\documentclass[11pt]{article}
\usepackage[margin=1in]{geometry}
\usepackage{amsmath, amssymb, amsthm}
\usepackage{graphicx, hyperref}
\usepackage{booktabs}
\usepackage{enumitem}
\usepackage{parskip}
\hypersetup{colorlinks=true, linkcolor=blue, urlcolor=blue, citecolor=blue}

\newtheorem{definition}{Definition}
\newtheorem{remark}{Remark}

\title{Statistics Primer for Whitespace 2}
\author{Originality Program --- Whitespace 2 Empirical Decomposition}
\date{April 2026}

\begin{document}
\maketitle
\tableofcontents
\clearpage

\section{Introduction}

This primer accompanies the Phase 0.1 scoping plan for Whitespace 2, an
empirical time-series study testing whether demographic plurality and
semantic plurality of scientific output have diverged over 1970--2024.
The plan commits to specific statistical methods, some of which are
standard in scientometrics but may be unfamiliar outside that literature.
This document explains each committed method from first principles:
what question it answers, its formal definition, the intuition behind
it, when it works and when it fails, and how Whitespace 2 uses it.

The organization follows the structure of the analysis pipeline:

\begin{itemize}[nosep]
  \item \textbf{Part I} --- diversity metrics used for the demographic,
    semantic, and canonical-concentration time series (Sections 2--6).
  \item \textbf{Part II} --- the four co-primary divergence tests
    committed in Phase 0.2 pre-registration (Sections 7--10).
  \item \textbf{Part III} --- break-point analysis for locating when any
    divergence begins (Sections 11--12).
  \item \textbf{Part IV} --- auxiliary methods used across multiple tests:
    bootstrap confidence intervals, multiple-comparisons correction,
    diachronic embedding alignment, and multicollinearity diagnostics
    (Sections 13--16).
\end{itemize}

Each section is self-contained. Read linearly for a full build-up, or
jump to specific methods as needed. A glossary of symbols is in the
appendix.

\clearpage
\part{Diversity Metrics}

\section{Shannon Entropy}
\label{sec:shannon}

\textbf{Question answered.} Given a categorical distribution (e.g.,
authors' countries, paper cluster assignments), how ``spread out''
is the mass across categories?

\textbf{Formal definition.} For a distribution with probabilities
$p_1, \ldots, p_K$ over $K$ categories:
\[
  H(p) = -\sum_{i=1}^{K} p_i \log p_i
\]
By convention $0 \log 0 = 0$. The base of the log determines the units:
base-2 yields bits, base-$e$ yields nats, base-10 yields hartleys.
Whitespace 2 uses natural log (nats); the choice affects units but not
conclusions.

\textbf{Intuition.} Shannon entropy is the expected ``surprise'' of a
random draw. If one category dominates ($p_1 \to 1$), drawing that
category is unsurprising and $H \to 0$. If all categories are equally
likely ($p_i = 1/K$), each draw is maximally surprising and
$H = \log K$.

The quantity $\exp(H)$ is the \textit{effective number of categories}
--- the number of equally-common categories that would produce the
same entropy. If papers are distributed over 50 cluster IDs but half
the mass is in 3 clusters, $\exp(H)$ might be around 8: roughly
equivalent to being spread over 8 equally-common clusters.

\textbf{Finite-sample bias.} The naïve ``plug-in'' estimator
$\hat H = -\sum \hat p_i \log \hat p_i$ systematically underestimates
the true entropy. The leading-order bias is
\[
  \mathbb{E}[\hat H] - H(p) \approx -\frac{K-1}{2N}
\]
where $N$ is the sample size. For $K = 100$ and $N = 1000$,
the bias is about $0.05$ nats---non-negligible. The \textit{Miller-Madow
bias correction} adds $(K-1)/(2N)$ to the plug-in estimate:
\[
  \hat H_{\text{MM}} = \hat H + \frac{K-1}{2N}
\]
Whitespace 2 uses this correction on all Shannon entropy estimates.

\textbf{When it fails.}
\begin{itemize}[nosep]
  \item With many rare categories ($K \sim N$), even Miller-Madow
    underestimates. Consider Chao-Shen or NSB estimators in that regime
    (not used in ws2 primary analyses).
  \item Requires a discrete distribution. For continuous features
    (career stage in years), discretize into bins first; bin choice
    affects the estimate.
\end{itemize}

\textbf{Use in ws2.} Per-dimension demographic diversity (gender,
country, institution type, prestige tier, career stage bin, training
institution) uses Shannon entropy as the primary metric. Cluster
entropy (Section~\ref{sec:cluster-entropy}) is Shannon entropy applied
to paper-embedding cluster assignments.

\section{Gini-Simpson Index}
\label{sec:gini-simpson}

\textbf{Question answered.} Given a categorical distribution, what is
the probability that two independently drawn samples fall into different
categories?

\textbf{Formal definition.} For probabilities $p_1, \ldots, p_K$:
\[
  \text{GS}(p) = 1 - \sum_{i=1}^{K} p_i^2
\]
The Simpson index $D = \sum p_i^2$ is the probability that two draws
match; Gini-Simpson is its complement.

\textbf{Intuition.} If a single category dominates, two random draws
likely match, $D \to 1$, and $\text{GS} \to 0$. If categories are
equal, the match probability is $1/K$ and $\text{GS} = 1 - 1/K \to 1$
as $K \to \infty$.

The quantity $1/D = 1/\sum p_i^2$ is also an \textit{effective number
of categories} (under a different weighting than Shannon's).
Sometimes called the ``Hill number of order 2.''

\textbf{How it differs from Shannon.} Both measure spread over
categories, but they weight rare vs. common categories differently:

\begin{itemize}[nosep]
  \item \textbf{Shannon} weights rare categories heavily (the
    $\log p_i$ term diverges as $p_i \to 0$). A rise in a few new rare
    categories moves Shannon entropy more.
  \item \textbf{Gini-Simpson} weights common categories heavily (the
    $p_i^2$ term). A rebalancing of existing dominant categories moves
    Gini-Simpson more.
\end{itemize}

A trend that is picked up by both suggests broad-based diversification.
A trend that moves only one suggests something more specific about the
shape of the change.

\textbf{Finite-sample bias.} Gini-Simpson's plug-in bias is smaller
than Shannon's:
\[
  \mathbb{E}[\widehat{\text{GS}}] - \text{GS}(p) \approx
  -\frac{1 - \text{GS}(p)}{N}
\]
Often acceptable without correction for $N \gtrsim 1000$.

\textbf{Use in ws2.} Secondary per-dimension demographic metric, paired
with Shannon entropy. Co-movement of both under trend testing is part
of the metric-plurality robustness commitment (desiderata §8).

\section{Rao's Quadratic Entropy}
\label{sec:rao-q}

\textbf{Question answered.} Given a categorical distribution
\emph{plus} a distance between categories, what is the expected pairwise
distance between two random samples from the distribution?

\textbf{Formal definition.} For probabilities $p_1, \ldots, p_K$ and a
symmetric distance function $d_{ij} \ge 0$ (with $d_{ii} = 0$):
\[
  Q(p, d) = \sum_{i=1}^{K} \sum_{j=1}^{K} p_i p_j d_{ij}
\]

\textbf{Intuition.} Rao's $Q$ generalizes Gini-Simpson. If $d_{ij} = 1$
for all $i \ne j$ (everything is equally distant), $Q = 1 - \sum p_i^2$
= Gini-Simpson. By specifying non-trivial distances between categories,
$Q$ captures the geometric sense in which the distribution is
``spread out'' in feature space.

\textbf{Distance choice is load-bearing.} The distance metric $d_{ij}$
is a modeling decision. For joint demographic states (e.g., female +
Chinese-affiliation + senior + non-top-50-institution), Whitespace 2
uses \textit{weighted Hamming distance}:
\[
  d_{ij} = \sum_{d=1}^{D} w_d \cdot \mathbf{1}(\text{state } i
  \text{ and } j \text{ differ on dimension } d)
\]
with uniform weights $w_d = 1/D$ in the primary analysis. Sensitivity
analysis doubles one weight at a time.

\textbf{Why not simpler composites?} Plain joint Shannon entropy over
joint states treats all category pairs as equally different. For ws2,
(female, American, senior) and (male, American, senior) should not be
\textit{as} far apart as (female, American, senior) and (male, Chinese,
early-career). Rao's $Q$ captures this; joint Shannon entropy does not.

\textbf{When it fails.}
\begin{itemize}[nosep]
  \item Requires a principled distance metric. If the metric is
    arbitrary, so is the conclusion.
  \item High-dimensional joint states may have sparse marginals,
    making $\hat p_i$ noisy. Bootstrap CIs are essential.
\end{itemize}

\textbf{Use in ws2.} Primary composite demographic-diversity metric.
Also used for paper-level \textit{team} demographic diversity in Test IV
(Section~\ref{sec:test-iv}).

\section{Effective Dimensionality (Participation Ratio)}
\label{sec:participation-ratio}

\textbf{Question answered.} Given a set of $N$ vectors in
$\mathbb{R}^d$ (e.g., paper embeddings), how many ``independent
directions'' does the cloud effectively span?

\textbf{Formal definition.} Let the empirical covariance matrix
of the $N$ vectors be $\Sigma = (1/N) \sum_{n=1}^N (x_n - \bar x)
(x_n - \bar x)^\top$, with eigenvalues $\lambda_1 \ge \lambda_2
\ge \cdots \ge \lambda_d \ge 0$. The \textit{participation ratio}
(also called effective rank or effective dimensionality) is:
\[
  d_{\text{eff}} = \frac{(\sum_i \lambda_i)^2}{\sum_i \lambda_i^2}
  = \frac{(\text{tr } \Sigma)^2}{\text{tr}(\Sigma^2)}
\]

\textbf{Intuition.} If all variance is in one direction
($\lambda_1 > 0$, rest zero), then $d_{\text{eff}} = 1$. If variance is
spread uniformly across $k$ directions
($\lambda_1 = \cdots = \lambda_k = \sigma^2$, rest zero), then
$d_{\text{eff}} = (k \sigma^2)^2 / (k \sigma^4) = k$. So
$d_{\text{eff}}$ counts ``how many roughly-equal eigenvalues there are,''
which corresponds to how many independent conceptual directions the
cloud inhabits.

\textbf{Why not just count eigenvalues above a threshold?} Hard cutoffs
are arbitrary and discontinuous. Participation ratio is smooth, scale-
invariant (multiplying all eigenvalues by a constant doesn't change it),
and has a clean probabilistic interpretation: it is the exponential of
the Rényi-2 entropy of the normalized eigenvalue distribution.

\textbf{Sample-size dependence.} For SPECTER2 embeddings
($d = 768$), estimating $d_{\text{eff}}$ requires $N \gtrsim d$ for a
stable estimate. The ratio itself converges faster than individual
eigenvalues, but small samples bias $d_{\text{eff}}$ downward.
Whitespace 2 calibrates the minimum sample size empirically in Phase 0.1
sanity Check 5b.

\textbf{Year-centering.} Embedding distributions often have a drifting
\textit{mean} (the centroid of CS 1975 sits in a different region than
CS 2024), which is a measurement artifact rather than a diversity signal.
Whitespace 2 subtracts the year-specific mean before computing
$d_{\text{eff}}$, so we measure \textit{spread around center}, not
absolute position.

\textbf{Use in ws2.} Primary semantic-diversity metric B (alongside
cluster entropy as primary A, mean pairwise cosine distance as
secondary).

\textbf{References.} Gao et al.~2017 (neuroscience); Recanatesi et
al.~2022 (scale-dependent dimensionality).

\section{Cluster Entropy}
\label{sec:cluster-entropy}

\textbf{Question answered.} Given embeddings clustered into $K$
groups, what is the Shannon entropy of each year's distribution over
clusters?

\textbf{Formal definition.} Let $c_n \in \{1, \ldots, K\}$ be the
cluster assignment of paper $n$. For year $Y$ with papers
$\mathcal{P}_Y$:
\[
  p_k^{(Y)} = \frac{|\{n \in \mathcal{P}_Y : c_n = k\}|}{|\mathcal{P}_Y|},
  \qquad
  H_Y = -\sum_{k=1}^{K} p_k^{(Y)} \log p_k^{(Y)}
\]
Apply Miller-Madow bias correction for finite $N$.

\textbf{Why cluster over the embedding space.} Embedding space is
continuous and high-dimensional; a direct Shannon entropy over the
continuous distribution requires density estimation, which is
expensive and sample-hungry. Clustering discretizes the space into
interpretable regions, after which Shannon entropy is straightforward.

\textbf{The critical methodological commitment: temporal stratification
of the cluster fit.} The clusters must be fit on a
\textit{temporally-stratified} subsample with equal papers per decade,
\textit{not} on the full pooled corpus.

If clusters are fit on the unstratified pool, modern papers (which
outnumber early-decade papers by 10--50$\times$) dominate the cluster
definitions. The clusters then reflect modern semantic structure; old
papers map poorly to those modern definitions and concentrate
artificially in a few least-wrong clusters, producing spurious
low-diversity in early years and a spurious divergence trend. This
would silently manufacture the paper's headline finding as a
corpus-composition artifact.

This is desiderata §11 and is load-bearing; it is stress-tested in
sanity Check 5b (Phase 0.1) with an A/B comparison quantifying the
artifact size on pilot data.

\textbf{Procedural separation.}
\begin{itemize}[nosep]
  \item \textbf{Fit:} on a temporally-stratified subsample (equal
    papers per decade).
  \item \textbf{Assign:} every paper in the full corpus (all years) to
    its nearest cluster centroid from the stratified fit.
\end{itemize}

\textbf{Robustness.} Report results for $K \in \{30, 50, 100\}$. If
the headline trend direction is consistent across $K$, the finding is
robust to cluster-count choice.

\textbf{Use in ws2.} Primary semantic-diversity metric A. Pairs
naturally with per-dimension demographic Shannon entropy, enabling
direct comparison on the same mathematical quantity
(effective-number-of-bins interpretation) in the main 3-panel figure.

\section{Citation Gini and Chu-Evans Spearman Top-N}
\label{sec:canonical}

Canonical concentration has two distinct senses:
\begin{enumerate}[nosep,label=(\arabic*)]
  \item \textbf{Distributional concentration}: are citations concentrated
    on a few papers?
  \item \textbf{Temporal persistence}: are the \textit{same} papers at
    the top year after year?
\end{enumerate}
The actuator-homogenization hypothesis is primarily about (2). Both
are captured by Whitespace 2's primary and secondary metrics.

\subsection{Citation Gini (secondary)}

The Gini coefficient on citations received in year $Y$:
\[
  G_Y = \frac{\sum_{i,j} |c_i - c_j|}{2 N^2 \bar c}
\]
where $c_i$ is the number of citations received by paper $i$ in year $Y$
and $\bar c$ is the mean. $G = 0$: perfect equality; $G = 1$: one paper
gets all citations.

\textbf{Why ``citations received in year $Y$'' (framing A):} this
captures the \textit{flow of attention} in year $Y$, i.e., what the
field is drawing on now. The alternative (``citations ever received by
papers published in year $Y$'') confounds paper age with canonicity.
Framing A is a pre-registered commitment and aligns with the
Chu-Evans methodology.

\subsection{Chu-Evans Spearman top-$N$ correlation (primary)}

For each year $Y$, rank papers by citations received in $Y$ (framing A
above). Define the top-$N$ list $T_Y$. Compute the Spearman rank
correlation between $T_Y$ and $T_{Y-\Delta}$:
\[
  \rho_Y = \text{spearman}(T_Y, T_{Y-\Delta})
\]
High $\rho_Y$ $\Rightarrow$ the same papers are at the top year over
year $\Rightarrow$ stable canon. Low $\rho_Y$ $\Rightarrow$ canon
turns over.

\textbf{Pre-registered parameters:} $N = 50$, $\Delta = 5$. Robustness
at $N \in \{100, 500\}$, $\Delta \in \{3, 10\}$.

\textbf{Why this metric is load-bearing.} The subfield mechanism test
--- ``canonical-concentrated subfields show more demographic-semantic
divergence'' --- uses the Chu-Evans Spearman as the canonical-
concentration score. This is the paper's single most important mechanism
analysis.

\subsection{Why CD-index is deliberately excluded}

Park, Leahey \& Funk's CD-index is contested (Petersen-Holst-Macher
critiques demonstrated citation-inflation artifacts and dataset
extraction bugs). Including CD-index as a primary metric would put
ws2 on the wrong side of an active methodological war. Used only for
Discussion-section engagement, not measurement.

\clearpage
\part{Divergence Tests}

Whitespace 2 commits to four co-primary divergence tests, each asking
a distinct substantive question. They are not redundant; they are
complementary.

\section{Test I: Standardized-Gap Trend}
\label{sec:test-i}

\textbf{Question answered.} Is the contemporaneous gap between
demographic and semantic diversity growing over time?

\textbf{Construction.} Standardize both time series over the analysis
window:
\[
  z^{\text{Dem}}_Y = \frac{\text{DemDiv}_Y - \bar{\text{DemDiv}}}
  {\sigma_{\text{DemDiv}}},
  \qquad
  z^{\text{Sem}}_Y = \frac{\text{SemDiv}_Y - \bar{\text{SemDiv}}}
  {\sigma_{\text{SemDiv}}}
\]
Define
\[
  \text{Gap}_Y = z^{\text{Dem}}_Y - z^{\text{Sem}}_Y
\]
Test:
\[
  \text{Gap}_Y = \alpha + \beta \cdot Y + \epsilon_Y,
  \quad H_0: \beta = 0,
  \quad H_1: \beta > 0 \text{ (one-sided)}
\]

\textbf{The autocorrelation problem.} Time series have serial
correlation: $\epsilon_Y$ and $\epsilon_{Y-1}$ are not independent.
Ordinary least squares (OLS) treats them as independent and
\textit{underestimates} standard errors, producing false positives.

\textbf{Newey-West HAC standard errors.} The
heteroscedasticity-and-autocorrelation-consistent (HAC) variance
estimator corrects for this:
\[
  \widehat{\text{Var}}_{\text{HAC}}(\hat \beta) =
  (X^\top X)^{-1} \hat{\Omega} (X^\top X)^{-1}
\]
where
\[
  \hat{\Omega} = \sum_{t=1}^T \hat \epsilon_t^2 x_t x_t^\top
  + \sum_{l=1}^{L} w_l \sum_{t=l+1}^{T} \hat \epsilon_t \hat \epsilon_{t-l}
  (x_t x_{t-l}^\top + x_{t-l} x_t^\top)
\]
The weights $w_l = 1 - l/(L+1)$ (Bartlett kernel) down-weight far
lags. The bandwidth $L$ is typically chosen as $L = \lfloor
4(T/100)^{2/9} \rfloor$ (Newey-West rule) or by bandwidth-selection
procedures.

\textbf{Intuition.} Newey-West is OLS with inflated standard errors
to account for the fact that correlated residuals give less
independent information than a count of data points suggests.

\textbf{Mann-Kendall as robustness.} A nonparametric alternative.
Compute the sign statistic:
\[
  S = \sum_{i < j} \text{sign}(\text{Gap}_{Y_j} - \text{Gap}_{Y_i})
\]
Under the null of no trend, $S$ is approximately normal with known
variance. No distributional assumption on residuals; robust to
outliers. Pre-whitening handles autocorrelation (Hamed \& Rao 1998).

\textbf{Use in ws2.} Test I primary on 6 metric pairings per field.
Bonferroni-corrected across the 6 pairings within each field.

\section{Test II: Confound-Controlled Gap Regression}
\label{sec:test-ii}

\textbf{Question answered.} Does the gap between demographic and
semantic diversity grow over time \textit{after controlling for known
confounds} (field growth, team-size evolution, publishing-norm shifts)?

\textbf{Construction.} Regress Gap on time and controls:
\[
  \text{Gap}_Y = \beta_0 + \sum_{k=1}^{K} \beta_k C_{k,Y} + \beta_t Y
  + \epsilon_Y
\]
where $C_k$ are the seven pre-registered controls: $\log N_{\text{papers}}$,
trimmed mean team size, median references per paper, arXiv fraction,
$\log$ distinct active authors, field-entry rate, and subfield
composition (vector). Newey-West HAC standard errors.

\[
  H_0: \beta_t = 0, \qquad H_1: \beta_t > 0 \text{ (one-sided)}
\]

\textbf{Why not regress SemDiv on DemDiv + time?} A prior draft of
Test II regressed $\text{SemDiv}_Y = \beta_0 + \beta_1 \text{DemDiv}_Y
+ \sum \beta_k C_k + \beta_t Y + \epsilon$. That specification has an
interpretability problem: if DemDiv itself trends strongly, including
$Y$ partials time out of DemDiv's effect, and $\beta_t$ can absorb the
divergence we want to detect. The Gap-reformulation avoids this by
treating divergence as a single quantity and asking whether it trends.

\textbf{Multicollinearity.} The seven controls are all upward-trending
over 50 years and are correlated. Variance Inflation Factor (VIF) is
reported per variable:
\[
  \text{VIF}_k = \frac{1}{1 - R_k^2}
\]
where $R_k^2$ is the $R^2$ from regressing $C_k$ on the other controls.
$\text{VIF} > 10$ is a standard flag. Whitespace 2 reports VIF
alongside coefficients. Mandatory sensitivity analysis: re-fit Test II
(i) with all controls, (ii) with each control dropped in turn, (iii)
with subfield composition toggled. $\beta_t$ stability across these
specifications is the robustness claim.

\textbf{Use in ws2.} Test II primary alongside Tests I, III, IV.
Directly addresses the ``does divergence survive controls?''
reviewer question.

\section{Test III: De-trended Cross-Correlation and Granger Causality}
\label{sec:test-iii}

\textbf{Question answered.} Does demographic diversity predict semantic
diversity at some lag $k$? If yes, the ``divergence'' may be
\textit{delayed tracking} rather than pure decoupling.

\subsection{Cross-correlation at lags}

The cross-correlation at lag $k$ between two series $X_t, Y_t$:
\[
  \rho_k = \frac{\text{Cov}(X_{t-k}, Y_t)}{\sqrt{\text{Var}(X_t) \cdot
  \text{Var}(Y_t)}}
\]

\textbf{The spurious-correlation trap.} If both $X_t$ and $Y_t$ have
strong trends (both rising over time), their raw cross-correlation
will be spuriously high --- not because they co-vary meaningfully, but
because both are going up. A textbook pitfall.

\textbf{Solution: de-trend first.} Two common approaches:
\begin{itemize}[nosep]
  \item \textbf{First differences:} $\Delta X_t = X_t - X_{t-1}$.
    Tests whether \textit{changes} in $X$ correlate with \textit{changes}
    in $Y$. Simple, aggressive de-trending.
  \item \textbf{Hodrick-Prescott (HP) filter:} extracts a smooth trend
    component and leaves cyclical residuals. Gentler; less aggressive.
\end{itemize}
Whitespace 2 uses first differences as primary de-trending and HP-filter
residuals as robustness.

\textbf{Pre-registered lag set.} $k \in \{0, 3, 5, 10, 15\}$. Fixed in
advance to limit multiple testing. Scanning all lags post-hoc and
reporting the largest inflates false-positive rate.

\textbf{Confidence bands.} Under the null hypothesis of independence,
$\rho_k$ is approximately normal with variance $1/T$ (Bartlett's
formula for sample sizes $T$). Confidence bands at $\pm 1.96/\sqrt{T}$.
Bootstrap CIs are a robustness option.

\subsection{Granger causality}

Formalizes ``does $X$ predict $Y$ beyond $Y$'s own history?''
Fit two models:
\begin{align*}
  \text{Restricted:}\quad Y_t &= \alpha + \sum_{l=1}^{L} \phi_l Y_{t-l}
  + u_t \\
  \text{Unrestricted:}\quad Y_t &= \alpha + \sum_{l=1}^{L} \phi_l Y_{t-l}
  + \sum_{l=1}^{L} \psi_l X_{t-l} + v_t
\end{align*}
F-test: is the improvement from adding $X$-lags significant?
\[
  F = \frac{(SSR_R - SSR_U)/L}{SSR_U/(T - 2L - 1)}
\]

\textbf{Caveats.}
\begin{itemize}[nosep]
  \item Granger causality is \textit{not} causality in the structural
    sense. It means $X$ has predictive information for $Y$ beyond $Y$'s
    own past. Confounders can manufacture or destroy Granger-causal
    relationships.
  \item Lag order $L$ matters. Chosen by AIC (Akaike information
    criterion) or BIC on pre-pilot data, capped at $L = 5$ in ws2.
  \item Requires stationarity of both series (after de-trending).
\end{itemize}

\textbf{Use in ws2.} Test III primary. Four-quadrant interpretation
combines Tests I and III:

\begin{center}
\begin{tabular}{lll}
  \toprule
  & \textbf{Test III: no lagged correlation} & \textbf{Test III: peak at $k > 0$} \\
  \midrule
  \textbf{Test I: significant} & true divergence & delayed tracking \\
  \textbf{Test I: not significant} & contemporaneous tracking & lagged tracking \\
  \bottomrule
\end{tabular}
\end{center}

Each quadrant is publishable and tells a different mechanism story.

\section{Test IV: Paper-Level Team-Diversity vs.\ Novelty}
\label{sec:test-iv}

\textbf{Question answered.} Do demographically diverse teams produce
semantically novel papers?

This is the within-paper companion to the aggregate Tests I--III.
It operates at a different level of aggregation (paper cross-section)
and addresses the intra-paper version of the claim: \textit{if} diverse
inputs get averaged toward canonical outputs by shared actuators, we
should see a within-team signature of it.

\subsection{Variables}

For each paper $p$:
\begin{itemize}[nosep]
  \item \textbf{Team diversity $T_p$}: Rao's $Q$ over the author team's
    joint demographic distribution, with uniform Hamming distance
    (same metric as the aggregate composite). $T_p = 0$ for single-author
    papers.
  \item \textbf{Novelty $N_p$ (primary)}: cosine distance from paper
    $p$'s embedding to the centroid of embeddings of papers $p$ cites:
    \[
      N_p = 1 - \frac{e_p \cdot \bar e_{\text{ref}(p)}}
      {\|e_p\| \cdot \|\bar e_{\text{ref}(p)}\|}
    \]
    High $N_p$ $\Rightarrow$ $p$ synthesizes distant concepts; low
    $N_p$ $\Rightarrow$ $p$ is close to its cited context.
  \item \textbf{Novelty $N_p$ (secondary)}: cosine distance from $p$
    to the year-$Y$ canonical centroid (top-100 most-cited papers in
    year $Y$). Tests ``distance from canon'' as alternative framing.
\end{itemize}

\subsection{Regression}

\[
  N_p = \gamma_0 + \gamma_1 T_p + \sum_k \gamma_k X_{k,p}
  + \mu_{\text{year}(p)} + \eta_{\text{subfield}(p)} + \epsilon_p
\]

where $X_{k,p}$ are paper-level controls (team size, reference count,
mean career stage, mean institutional prestige, log paper age, field
dummy) and $\mu, \eta$ are year and subfield fixed effects.

\textbf{Fixed effects, intuitively.} A fixed effect absorbs all
variation due to year or subfield, leaving only variation
\textit{within} year-subfield cells. This answers ``within a given year
and subfield, do diverse teams produce more or less novel papers?''
Year-trends and subfield-differences don't confound the estimate.

\subsection{Standard errors: double-clustered}

Papers by the same lead author or in the same subfield are not
independent. Double-clustering the standard errors accounts for this:
\[
  \widehat{\text{Var}}_{\text{DC}}(\hat \gamma_1) =
  V_{\text{author}} + V_{\text{subfield}} - V_{\text{author} \cap
  \text{subfield}}
\]
(Cameron, Gelbach \& Miller 2011 estimator). Naïve OLS standard errors
that ignore clustering underestimate variance and produce false
positives.

\subsection{Hypotheses and effect size}

\begin{itemize}[nosep]
  \item $H_0$: $\gamma_1 = 0$ (team diversity uncorrelated with novelty)
  \item $H_1$ (actuator-dominant): $\gamma_1 < 0$ --- diverse teams
    produce \textit{less} novel work; shared actuators average diverse
    inputs toward canonical outputs.
  \item $H_1'$ (diversity-productive): $\gamma_1 > 0$ --- diverse teams
    produce more novel work. Consistent with Hofstra et al.~2020
    framing.
\end{itemize}

Two-sided test; the \textit{sign} of $\gamma_1$ is the substantive
answer. With hundreds of thousands of papers, statistical significance
is almost automatic, so substantive-significance threshold
$|\gamma_1| \ge 0.05$ (after standardization of $T_p$ and $N_p$) is
pre-registered.

\textbf{Time interaction (secondary).}
$\gamma_1(Y) = \gamma_{10} + \gamma_{11} Y$: does the relationship
evolve? A strengthening negative $\gamma_1$ over time would be
consistent with growing actuator dominance.

\textbf{Causal caveat.} Cross-sectional correlations $\ne$ causal
effects. $\gamma_1$ is reported as descriptive association. Reverse
causation (novel work attracts diverse teams), selection (risk-prefering
researchers form diverse teams), and omitted variables all threaten
causal interpretation. No instrumental variable strategy is attempted
in Whitespace 2.

\clearpage
\part{Break-Point Analysis}

If Tests I or II detect significant divergence, a natural follow-up is:
when does it start? A continuous slow divergence from 1970 tells a
different mechanism story than a sharp break around, say, 2000.

\section{Chow Test at Pre-specified Candidates}
\label{sec:chow}

\textbf{Setup.} Suppose a candidate break year $\tau$ is specified in
advance (e.g., 2000, corresponding to the dot-com peak). Fit two
regressions of Gap on year:
\begin{align*}
  \text{Pre:}& \quad \text{Gap}_Y = \alpha_1 + \beta_1 Y + u_Y,
  \quad Y \le \tau \\
  \text{Post:}& \quad \text{Gap}_Y = \alpha_2 + \beta_2 Y + v_Y,
  \quad Y > \tau \\
  \text{Pooled:}& \quad \text{Gap}_Y = \alpha + \beta Y + w_Y,
  \quad \text{all } Y
\end{align*}
Let $SSR_P, SSR_1, SSR_2$ be the residual sums of squares. The Chow
F-statistic:
\[
  F = \frac{(SSR_P - SSR_1 - SSR_2)/2}{(SSR_1 + SSR_2)/(T-4)}
\]
tests the null that coefficients are equal on both sides of $\tau$.

\textbf{Pre-specification is the key.} Post-hoc search for break
points inflates false positives enormously --- with 50 years of data,
\textit{some} year will look like a break by chance. Chow requires
the candidate to be committed in advance. Whitespace 2 pre-registers
specific candidates per field, tied to substantive events.

\textbf{Pre-registered candidates (CS):}
\begin{itemize}[nosep]
  \item 1991--93: arXiv launch + public web
  \item 1998--2000: dot-com peak (demographic broadening from industry
    demand)
  \item 2008--09: financial crisis, academic-market consolidation
  \item 2012: AlexNet, deep-learning takeoff
  \item 2018--20: foundation models (likely too recent for smoothed
    series)
\end{itemize}

\textbf{Pre-registered candidates (Physics):}
\begin{itemize}[nosep]
  \item 1991: arXiv launched in hep-th
  \item 1995--2000: post-SSC cancellation; string theory dominance
  \item 2012: Higgs boson discovery
\end{itemize}

Bonferroni correction across the candidate set within each field.

\section{Bai-Perron Multi-Break Detection}
\label{sec:bai-perron}

\textbf{Setup.} Rather than pre-specifying candidates, the Bai-Perron
procedure algorithmically searches for $m$ breakpoints
$\tau_1 < \tau_2 < \cdots < \tau_m$ partitioning the time series into
$m+1$ regimes, each with its own slope and intercept. The number of
breaks $m$ is selected by BIC or sequential F-tests.

\textbf{Trimming.} Breaks cannot be detected near series boundaries
(too few points on one side for stable estimates). Standard trimming
removes the first and last 15\% of the series. For 1970--2024, this
excludes break detection before $\sim$1978 or after $\sim$2017.

\textbf{Confidence intervals on break locations.} Bai-Perron provides
asymptotic CIs on each $\tau_i$. Bootstrap CIs are a robustness option.

\textbf{Power.} With annual data over 55 years, the Bai-Perron
procedure can reliably detect 0 or 1 breaks; 2 breaks is borderline;
3+ is generally not defensible at this series length. Whitespace 2
reports what BIC selects plus sensitivity at $\pm 1$ break.

\textbf{Concordance requirement.} A break is ``significant'' only if
\textit{both} Bai-Perron detects it \textit{and} Chow at a nearby
pre-specified candidate (within $\pm 2$ years) rejects no-break.
Prevents one-method-only artifacts from being reported as substantive.

\clearpage
\part{Auxiliary Methods}

\section{Bootstrap Confidence Intervals}
\label{sec:bootstrap}

\textbf{What and why.} Classical confidence intervals rely on
distributional assumptions (usually normality) about the sampling
distribution of the statistic. For complex statistics like Rao's $Q$
or participation-ratio effective dimensionality, these assumptions
may not hold. The bootstrap sidesteps the issue.

\textbf{Algorithm.} Given data $\mathcal{D} = \{x_1, \ldots, x_N\}$ and
statistic $\hat\theta(\mathcal{D})$:
\begin{enumerate}[nosep]
  \item For $b = 1, \ldots, B$: draw a sample $\mathcal{D}^{*(b)}$ of
    size $N$ with replacement from $\mathcal{D}$; compute
    $\hat\theta^{*(b)} = \hat\theta(\mathcal{D}^{*(b)})$.
  \item Form the empirical distribution of
    $\{\hat\theta^{*(1)}, \ldots, \hat\theta^{*(B)}\}$.
  \item For 95\% CI: take the 2.5th and 97.5th percentiles (percentile
    method) or the bias-corrected equivalent (BCa method).
\end{enumerate}

\textbf{Commitment.} Whitespace 2 uses $B = 200$ replicates for every
metric on every figure. 95\% percentile CIs unless otherwise noted.

\textbf{When it fails.}
\begin{itemize}[nosep]
  \item Small sample sizes ($N < 30$) give noisy bootstrap estimates.
  \item Statistics at distributional extremes (maxima, minima) are
    poorly bootstrapped. Not an issue for our metrics.
  \item Time-series bootstrap requires block bootstrapping
    (Politis-Romano) to preserve autocorrelation structure. For our
    annual series in the divergence tests, we use moving-block
    bootstrap when computing CIs on trend estimates.
\end{itemize}

\section{Multiple-Comparisons Correction}
\label{sec:multiple-comparisons}

\textbf{The problem.} Running many statistical tests inflates the
probability of at least one false positive. With 100 tests at
$\alpha = 0.05$ under the global null, we expect 5 false positives
just by chance.

\subsection{Bonferroni correction}

Divide $\alpha$ by the number of tests:
\[
  \alpha_{\text{Bonf}} = \frac{\alpha}{m}
\]
Controls the \textit{family-wise error rate} (FWER): probability of
\textit{any} false positive among the $m$ tests. Conservative; loses
power when tests are correlated (since true effective $m$ is smaller).

\subsection{False discovery rate (FDR)}

Benjamini-Hochberg procedure. Sort p-values
$p_{(1)} \le \cdots \le p_{(m)}$. Reject tests $p_{(i)}$ where
$p_{(i)} \le (i/m) \cdot q$ for target FDR level $q$ (often 0.05 or 0.10).

Controls the \textit{expected proportion} of false discoveries among
rejections. Less conservative than Bonferroni; more powerful when
many true effects exist.

\subsection{Whitespace 2's hierarchical scheme}

\begin{itemize}[nosep]
  \item Primary tests (Tests I--IV, 6 metric pairings each, 2 fields):
    Bonferroni within each test-type across 6 pairings;
    \textit{no correction across test-types} since they ask distinct
    questions. Replicated divergence requires all 3 test types to agree
    on direction in $\ge$4/6 pairings per field.
  \item Secondary tests (per-dimension demographic decomposition:
    7 dims $\times$ 3 semantic $\times$ 2 fields = 42 tests): FDR at
    $q = 0.10$, labeled exploratory.
  \item Break-point analysis: Bonferroni within candidate set per field.
\end{itemize}

\section{Procrustes Alignment}
\label{sec:procrustes}

\textbf{Problem setup.} In the Flavor A drift mitigation, we train
separate embedding spaces per era and then need to \textit{align} them
so embeddings from different eras can be compared. Each space is defined
up to an arbitrary rotation/reflection.

\textbf{Orthogonal Procrustes problem.} Given two matrices
$W_1, W_2 \in \mathbb{R}^{k \times d}$ (each row a vector), find the
orthogonal matrix $R$ minimizing:
\[
  R^* = \arg\min_{R : R^\top R = I} \| W_1 R - W_2 \|_F
\]

\textbf{Closed-form solution via SVD.} Compute
$W_2^\top W_1 = U \Sigma V^\top$. Then $R^* = V U^\top$.

\textbf{The anchor-paper idea.} To align era 1 (say, 1980s embeddings)
with era 2 (2010s embeddings), we need vectors that mean the ``same
thing'' in both eras. These are anchor papers or anchor concepts
appearing in both eras. Their embeddings in era 1 form $W_1$; their
embeddings in era 2 form $W_2$. Procrustes finds the rotation aligning
these anchors. Then apply the same rotation to all era-1 embeddings to
bring them into era-2's coordinate system.

\textbf{Anchor-choice caveat.} Whatever makes it into the anchor set
defines the ``fixed frame.'' Anchors should be semantically stable
across eras (named methods, classical theorems, widely-agreed terms).
This is itself a modeling commitment; Whitespace 2 treats anchor
selection as a Phase 0.2 deliverable.

\section{Variance Inflation Factor (VIF)}
\label{sec:vif}

\textbf{Problem.} In regression with multiple controls, if controls
are highly correlated, coefficient estimates become unstable and their
standard errors explode. This is multicollinearity.

\textbf{Diagnostic.} For control $k$, regress it on the other controls
and record $R_k^2$. Define:
\[
  \text{VIF}_k = \frac{1}{1 - R_k^2}
\]
$\text{VIF}_k = 1$: control $k$ is uncorrelated with others.
$\text{VIF}_k = 10$: $R_k^2 = 0.9$; 90\% of $k$'s variance is
explained by other controls. Standard error of $\hat\beta_k$ is
inflated by a factor $\sqrt{\text{VIF}_k}$.

\textbf{Rules of thumb.}
\begin{itemize}[nosep]
  \item VIF $< 5$: unproblematic.
  \item VIF 5--10: concerning; re-examine.
  \item VIF $> 10$: drop, combine, or report stability under omission.
\end{itemize}

\textbf{Whitespace 2 treatment.} The seven Test II controls are all
upward-trending over 50 years; high VIFs are \textit{expected}. Rather
than dropping variables mechanically, ws2 commits to a
\textit{sensitivity analysis}: re-fit with each control dropped in turn,
report $\beta_t$ across specifications. Qualitative stability of
$\beta_t$ is the robustness claim.

\section{Missing Data: MAR, MNAR, and Inverse-Probability Weighting}
\label{sec:mar-mnar}

\textbf{Problem.} Whitespace 2's analytical population $P$ — papers
with non-empty \texttt{abstract\_inverted\_index} in OpenAlex — is
$\sim$50\% of the OpenAlex CS+Physics field. Aggregate metrics computed
on $P$ are biased estimates of the field-level aggregates if the
selection mechanism (whether a paper has an abstract or not) correlates
with the metric being measured. Phase 0.1 Check 1f confirmed this is the
case along three axes: citation count, concept distribution, and
country of first authorship — all aligned with ws2's three substantive
measurement axes.

\textbf{Rubin's missing-data taxonomy} (Rubin 1976; standard in
econometrics, biostatistics, and survey statistics). Let $D$ be the
indicator that a paper has an abstract (1 if observed, 0 if missing),
let $X$ be the observed covariates (year, field, type, citation tertile,
concept cluster, known country), and let $Y$ be the outcome of interest
(e.g., paper-level semantic novelty, demographic attribute, citation
rank).

\begin{itemize}[nosep]
  \item \textbf{MCAR (Missing Completely At Random):}
        $\Pr(D = 1 \mid X, Y) = \Pr(D = 1)$. Missingness is independent
        of everything observable and unobservable. Aggregate means are
        unbiased without correction. \textit{ws2 does not assume MCAR};
        Check 1f rules it out empirically.
  \item \textbf{MAR (Missing At Random):}
        $\Pr(D = 1 \mid X, Y) = \Pr(D = 1 \mid X)$. Missingness depends
        on \emph{observed} covariates but not on the missing value
        itself, conditional on those covariates. Two papers with
        identical $X$ have the same probability of being observed
        regardless of what their outcome $Y$ would have been.
        \textbf{Inverse-probability-of-availability weighting recovers
        full-population aggregates exactly under MAR.}
  \item \textbf{MNAR (Missing Not At Random):}
        $\Pr(D = 1 \mid X, Y) \neq \Pr(D = 1 \mid X)$. Missingness
        depends on the unobserved $Y$ itself, even after conditioning
        on $X$. Recovered aggregates are biased; we can only
        \emph{bound} the population aggregate, not point-identify it.
\end{itemize}

\textbf{Inverse-probability weighting under MAR.} For each paper
$p \in P$, fit a propensity model
$\hat\pi(X_p) = \Pr(D = 1 \mid X = X_p)$ and define the weight
$w_p = 1 / \hat\pi(X_p)$. The weighted aggregate
\[
  \hat M_{\text{corrected}}
  = \frac{\sum_{p \in P} w_p \, m_p}{\sum_{p \in P} w_p}
\]
is an unbiased estimate of the OpenAlex full-population aggregate
$M$ under MAR. Intuition: papers with low $\hat\pi(X_p)$ are
underrepresented in $P$ (papers like them are rarely observed); their
weights are correspondingly large to compensate. Under MAR, the
conditional distribution of $Y$ given $X$ is the same in $P$ as in the
full population, so re-weighting by $1/\hat\pi$ recovers the population
distribution.

\textbf{Sensitivity to MNAR.} The MAR assumption is testable in
principle but never directly verifiable (we cannot observe $Y$ for
$D = 0$ papers). The standard approach is a \emph{sensitivity grid}:
parametrize a tilt $\gamma \in [-\Gamma, \Gamma]$ representing how much
the missingness mechanism deviates from MAR, recompute the bound
$\hat M(\gamma)$ for a range of $\gamma$, and report the bound
$[\hat M(-\Gamma), \hat M(+\Gamma)]$. ws2 commits to a pre-registered
$\Gamma$-grid in Phase 0.2; tentative $\Gamma$ corresponds to $\pm 50\%$
selection-probability variation within $X$-strata.

\textbf{Reliability bound on the propensity model.} The
inverse-probability-weighted estimate is only as good as the
propensity model. Pre-registered diagnostics:
\begin{itemize}[nosep]
  \item \textbf{Held-out CV-AUC} of the fitted propensity model.
        ws2 commits to a pre-registered threshold (tentative $\geq 0.75$
        in Phase 0.2). Below threshold, the corrected aggregate is not
        the headline; uncorrected + MNAR-tilt bound is reported.
  \item \textbf{Calibration plot} of predicted vs.\ observed $\hat\pi$
        across deciles. Severe miscalibration indicates the propensity
        model is misspecified; correction may be unreliable even at
        adequate AUC.
  \item \textbf{Stabilized weights} $w_p / \bar w$ to flag
        high-leverage observations (papers with extreme weights drive
        the aggregate; their over-influence inflates variance).
  \item \textbf{Functional-form sensitivity}: report aggregates under
        logistic regression (baseline) and gradient-boosted trees
        (flexibility check). Substantial disagreement is itself a
        diagnostic of nonlinear interactions the baseline misses.
\end{itemize}

\textbf{Whitespace 2 treatment.} \S9e in the Phase 0.1 plan
operationalizes inverse-probability weighting on
$X = \text{year} \times \text{field} \times \text{type} \times
\text{citation\_tertile} \times \text{concept\_cluster} \times
\text{known\_country}$. Headline aggregates are reported with both
the MAR-corrected estimate (when CV-AUC clears the threshold) and
the MNAR-tilt bound (always). The country axis is not propensity-
correctable (the UNKNOWN bucket has no observable strata to fit on);
demographic-plurality claims are scope-narrowed to $P_{\text{demo}}$
($\sim$50\% with determinable country) by construction. See plan
\S9e for the full operational specification and \S0 for the
analytical-population definitions.

\textbf{Why MAR/MNAR is the right framework here.} Selection bias
caused by missing data is structurally distinct from
attribute-classification error (e.g., NamSor misgendering authors)
and from metric-construct invalidity (e.g., CD-index zero-reference
artifacts). The three are orthogonal construct-validity threats;
ws2's Methods overview (\S14) frames them as four parallel audits.
The MAR/MNAR taxonomy is the principled toolkit for the
analytical-population audit; without it, ``selection bias'' is too
fuzzy to bound and too tempting to wave away.

\clearpage
\part*{Appendix}

\section*{Glossary of Symbols}

\begin{tabular}{ll}
  \toprule
  Symbol & Meaning \\
  \midrule
  $H(p)$ & Shannon entropy of distribution $p$ \\
  $\text{GS}(p)$ & Gini-Simpson index \\
  $Q(p, d)$ & Rao's quadratic entropy with distance $d$ \\
  $d_{\text{eff}}$ & Effective dimensionality (participation ratio) \\
  $\lambda_i$ & $i$-th eigenvalue of covariance matrix \\
  $p_k^{(Y)}$ & Cluster-$k$ probability in year $Y$ \\
  $G_Y$ & Citation Gini in year $Y$ \\
  $\rho_Y$ & Chu-Evans Spearman top-$N$ at year $Y$, lag $\Delta$ \\
  $z^{\text{Dem}}_Y, z^{\text{Sem}}_Y$ & Standardized demographic, semantic diversity \\
  $\text{Gap}_Y$ & $z^{\text{Dem}}_Y - z^{\text{Sem}}_Y$ \\
  $\beta_t$ & Slope on year in Test II \\
  $T_p$ & Team demographic diversity of paper $p$ \\
  $N_p$ & Novelty of paper $p$ \\
  $\gamma_1$ & Coefficient of $T_p$ in Test IV regression \\
  $\mu, \eta$ & Year and subfield fixed effects \\
  $\tau$ & Break-point year \\
  $\text{VIF}_k$ & Variance inflation factor for control $k$ \\
  $B$ & Bootstrap replicates \\
  \bottomrule
\end{tabular}

\section*{Key References by Method}

\begin{itemize}[nosep]
  \item \textbf{Shannon entropy, Miller-Madow correction:} Miller 1955
    ``Note on the bias of information estimates.''
  \item \textbf{Rao's Q:} Rao 1982, ``Diversity and dissimilarity
    coefficients: a unified approach,'' \textit{Theor.\ Popul.\ Biol.}
  \item \textbf{Participation ratio:} Gao et al.\ 2017 (neural
    dimensionality); Recanatesi et al.\ 2022.
  \item \textbf{Chu-Evans Spearman:} Chu \& Evans 2021,
    \textit{PNAS}~118(41).
  \item \textbf{Newey-West HAC:} Newey \& West 1987,
    \textit{Econometrica}.
  \item \textbf{Mann-Kendall with pre-whitening:} Hamed \& Rao 1998,
    \textit{J.~Hydrology}.
  \item \textbf{Granger causality:} Granger 1969,
    \textit{Econometrica}.
  \item \textbf{Bai-Perron:} Bai \& Perron 1998, \textit{Econometrica}
    66(1); 2003, \textit{J.~Appl.\ Econ.}
  \item \textbf{Cameron-Gelbach-Miller double-clustering:}
    Cameron, Gelbach \& Miller 2011, \textit{J.~Bus.\ Econ.\ Stat.}
  \item \textbf{Bootstrap:} Efron \& Tibshirani 1993,
    \textit{An Introduction to the Bootstrap}.
  \item \textbf{Block bootstrap for time series:} Politis \& Romano 1994,
    \textit{JASA}.
  \item \textbf{FDR:} Benjamini \& Hochberg 1995, \textit{J.\ R.~Stat.\
    Soc.\ B}.
  \item \textbf{Orthogonal Procrustes:} Schönemann 1966,
    \textit{Psychometrika}.
  \item \textbf{Diachronic word embeddings:} Hamilton, Leskovec \&
    Jurafsky 2016, \textit{ACL}.
\end{itemize}

\section*{Building This Document}

From the primer directory:
\begin{verbatim}
pdflatex stats.tex
pdflatex stats.tex    # second pass for TOC
\end{verbatim}

No bibliography manager is used; references are in-text. If the
references grow, migrate to \texttt{bibtex} or \texttt{biblatex}.

\end{document}
